% ============================================================
\part{Linux praktisch verstehen}
% ============================================================

\chapter{Wir wollen eine Datei speichern und wiederfinden}

\kapitelziel{
Dieses Kapitel wird nach dem Werkstatt-Schema ausgearbeitet.
}

\kapitelstruktur

\begin{aufgabeBox}
Aufgabe folgt.
\end{aufgabeBox}

% TODO: Situation, Problem, Wissen, Lösung, Merksatz,
%       Fehler, Übungen, Quiz, Haube, Robotik-Transfer

\chapter{Wir wollen ein Programm starten}

\kapitelziel{
Dieses Kapitel wird nach dem Werkstatt-Schema ausgearbeitet.
}

\kapitelstruktur

\begin{aufgabeBox}
Aufgabe folgt.
\end{aufgabeBox}

% TODO: Situation, Problem, Wissen, Lösung, Merksatz,
%       Fehler, Übungen, Quiz, Haube, Robotik-Transfer

\chapter{Wir wollen Dateien erstellen, kopieren, verschieben und löschen}

\kapitelziel{
Dieses Kapitel wird nach dem Werkstatt-Schema ausgearbeitet.
}

\kapitelstruktur

\begin{aufgabeBox}
Aufgabe folgt.
\end{aufgabeBox}

% TODO: Situation, Problem, Wissen, Lösung, Merksatz,
%       Fehler, Übungen, Quiz, Haube, Robotik-Transfer

\chapter{Wir wollen wissen, was auf dem Computer gerade läuft}

\kapitelziel{
Dieses Kapitel wird nach dem Werkstatt-Schema ausgearbeitet.
}

\kapitelstruktur

\begin{aufgabeBox}
Aufgabe folgt.
\end{aufgabeBox}

% TODO: Situation, Problem, Wissen, Lösung, Merksatz,
%       Fehler, Übungen, Quiz, Haube, Robotik-Transfer

\chapter{Wir wollen Fehlermeldungen verstehen}

\kapitelziel{
Dieses Kapitel wird nach dem Werkstatt-Schema ausgearbeitet.
}

\kapitelstruktur

\begin{aufgabeBox}
Aufgabe folgt.
\end{aufgabeBox}

% TODO: Situation, Problem, Wissen, Lösung, Merksatz,
%       Fehler, Übungen, Quiz, Haube, Robotik-Transfer

